\documentclass[11pt]{article}

\usepackage[margin=1in]{geometry}
\usepackage[T1]{fontenc}
\usepackage{lmodern}
\usepackage{microtype}
\usepackage{amsmath,amssymb}
\usepackage{booktabs}
\usepackage{longtable}
\usepackage{tabularx}
\usepackage{array}
\usepackage{xcolor}
\usepackage{enumitem}
\usepackage{listings}
\usepackage[hidelinks]{hyperref}

\definecolor{codeblue}{RGB}{25,80,135}
\definecolor{codegreen}{RGB}{30,120,65}
\definecolor{codegray}{RGB}{90,90,90}
\definecolor{codebackground}{RGB}{247,247,247}

\lstdefinelanguage{json}{
  basicstyle=\ttfamily\small,
  numbers=left,
  numberstyle=\tiny\color{codegray},
  stepnumber=1,
  numbersep=8pt,
  showstringspaces=false,
  breaklines=true,
  frame=single,
  backgroundcolor=\color{codebackground},
  string=[s]{\"}{\"},
  stringstyle=\color{codegreen},
  literate=
   *{0}{{{\color{codeblue}0}}}{1}
    {1}{{{\color{codeblue}1}}}{1}
    {2}{{{\color{codeblue}2}}}{1}
    {3}{{{\color{codeblue}3}}}{1}
    {4}{{{\color{codeblue}4}}}{1}
    {5}{{{\color{codeblue}5}}}{1}
    {6}{{{\color{codeblue}6}}}{1}
    {7}{{{\color{codeblue}7}}}{1}
    {8}{{{\color{codeblue}8}}}{1}
    {9}{{{\color{codeblue}9}}}{1}
}

\lstset{
  basicstyle=\ttfamily\small,
  numbers=left,
  numberstyle=\tiny\color{codegray},
  numbersep=8pt,
  showstringspaces=false,
  breaklines=true,
  frame=single,
  backgroundcolor=\color{codebackground},
  keywordstyle=\color{codeblue}\bfseries,
  commentstyle=\color{codegray}\itshape,
  stringstyle=\color{codegreen}
}

\newcommand{\field}[1]{\texttt{#1}}
\newcommand{\vect}[1]{\boldsymbol{#1}}
\newcommand{\R}{\mathbb{R}}

\title{MDI Chimera Final-Patch JSON Interchange Format\\
       \large Version 1: Reference and Area-Computation Guide}
\author{MDI Chimera Project}
\date{\today}

\begin{document}
\maketitle
\tableofcontents
\newpage

\section{Purpose and scope}

This document describes version~1 of the JSON interchange format produced by
MDI Chimera.  The format is intended to let an independent program reproduce
the final intersection patches between a Cartesian volume grid and a spherical
surface grid.  In particular, it is designed to support independent computation
of patch areas.

The file is not a serialization of Java implementation objects.  It is a
language-neutral geometric description containing:

\begin{itemize}[nosep]
  \item the coordinate conventions and units;
  \item the complete Cartesian and spherical grids;
  \item one zero-based five-tuple of grid indices for every final patch;
  \item an ordered, directed, closed boundary for every patch;
  \item exact parametric descriptions of all boundary-curve types; and
  \item numerical diagnostics computed by Chimera for comparison.
\end{itemize}

The authoritative machine-readable definition is
\field{mdi-chimera-patches-v1.schema.json}, distributed with Chimera.  A reader
should check both the root \,\field{format}\, and \,\field{version}\, fields before
interpreting the remainder of the document.

\section{Coordinate conventions}

Version~1 currently uses the GSM coordinate system and Earth radii as its
length unit, although both values are stated explicitly in the file.  All
three-component arrays are ordered
\[
  [x,y,z].
\]
All angular quantities are in radians.

On a sphere of radius $R$ centered at the origin, Chimera uses colatitude
$\theta$ and azimuth $\phi$:
\begin{align}
  x &= R\sin\theta\cos\phi,\\
  y &= R\sin\theta\sin\phi,\\
  z &= R\cos\theta.
\end{align}
Thus $\theta=0$ is the positive-$z$ (north) pole, $\theta=\pi/2$ is the
equator, and $\theta=\pi$ is the negative-$z$ (south) pole.  The azimuth $\phi$
is measured in the $x$--$y$ plane from positive $x$ toward positive $y$ and is
normally reported in $[-\pi,\pi]$.

Every curve has a parameter $t\in[0,1]$.  A curve is \emph{directed}: $t=0$
is its start and $t=1$ is its end.  Signed angular increments must never be
replaced blindly by absolute values or normalized to a short arc.

\section{Top-level organization}

An abbreviated file has the following form.

\begin{lstlisting}[language=json,caption={Abbreviated version-1 document}]
{
  "format" : "mdi-chimera-patches",
  "version" : 1,
  "conventions" : { ... },
  "grid" : { ... },
  "summary" : { ... },
  "patches" : [
    {
      "id" : 0,
      "indices" : {
        "nx" : 1, "ny" : 2, "nz" : 3,
        "nTheta" : 4, "nPhi" : 5
      },
      "diagnostics" : { ... },
      "boundary" : [ ... directed curves ... ]
    }
  ]
}
\end{lstlisting}

\subsection{Root fields}

\begin{longtable}{>{\ttfamily}p{0.21\textwidth}p{0.13\textwidth}p{0.58\textwidth}}
\toprule
\normalfont Field & Type & Meaning \\
\midrule
\endhead
format & string & Must equal \field{mdi-chimera-patches}. \\
version & integer & Must equal $1$ for the format described here. \\
conventions & object & Units, angular definitions, and traversal conventions. \\
grid & object & Complete input-grid definition. \\
summary & object & File-wide counts and numerical diagnostics. \\
patches & array & Final patches.  Array order determines the exported IDs. \\
\bottomrule
\end{longtable}

\subsection{The conventions object}

\begin{longtable}{>{\ttfamily}p{0.27\textwidth}p{0.65\textwidth}}
\toprule
\normalfont Field & Version-1 interpretation \\
\midrule
\endhead
lengthUnit & Currently \field{EARTH\_RADII}.  All Cartesian lengths, radii,
and perimeters use this unit. \\
coordinateSystem & Currently \field{GSM}. \\
angleUnit & \field{RADIANS}. \\
theta & \field{COLATITUDE\_FROM\_POSITIVE\_Z}. \\
phi & \field{AZIMUTH\_FROM\_POSITIVE\_X\_TOWARD\_POSITIVE\_Y}. \\
phiRange & Normally $[-\pi,\pi]$. \\
curveParameterRange & $[0,1]$. \\
boundaryOrientation & \field{ORDERED\_DIRECTED\_CLOSED\_LOOP}. \\
\bottomrule
\end{longtable}

\section{Grid definition}

\subsection{Sphere}

The \,\field{grid.sphere}\, object contains a three-vector \,\field{center}\,
and scalar \,\field{radius}.  Version~1 readers should require the center to be
the origin.  The radius $R$ is used in every curve and area formula.

\subsection{Cartesian grid}

The arrays \,\field{xVertices}, \,\field{yVertices}, and \,\field{zVertices}
contain strictly increasing local grid coordinates.  If
\[
  \vect{o}=(o_x,o_y,o_z)=\field{offset},
\]
then the global coordinate of a vertex is its stored coordinate plus the
corresponding offset.  For example, Cartesian cell $n_x=i$ occupies
\[
  x_i+o_x \le x < x_{i+1}+o_x,
\]
with the usual inclusion of the final upper boundary.  The $y$ and $z$ axes
are analogous.  With $N_x$ vertices there are $N_x-1$ cells and valid indices
are $0,\ldots,N_x-2$.

\subsection{Spherical grid}

The arrays \,\field{thetaVertices}\, and \,\field{phiVertices}\, are strictly
increasing.  The theta array spans $[0,\pi]$ and the phi array spans
$[-\pi,\pi]$.  Spherical cell $(n_\theta,n_\phi)=(j,k)$ is bounded by
\begin{align}
  \theta_j &\le \theta < \theta_{j+1},\\
  \phi_k &\le \phi < \phi_{k+1},
\end{align}
again treating the final upper endpoint as part of the last cell.

\section{Patches and the five-tuple}

Each patch contains the indices
\[
  (n_x,n_y,n_z,n_\theta,n_\phi).
\]
They are zero-based indices into the four grid-vertex arrays.  The five-tuple
identifies the Cartesian volume cell and spherical surface cell whose
intersection is the patch.  Within a valid completed Chimera result, five-tuples
are unique.

The \,\field{boundary}\, array is not an unordered set.  If it contains curves
$C_0,\ldots,C_{m-1}$, they must be traversed in that order, and
\[
  C_i(1)=C_{i+1}(0),\qquad
  C_{m-1}(1)=C_0(0),
\]
up to floating-point tolerance.  Changing order, independently reversing
curves, or discarding signs destroys the boundary topology.

The integer \,\field{id}\, is a convenient file-local identifier.  It is not a
geometric invariant and should not be expected to remain stable if patches are
regenerated with a different Chimera version or ordering rule.

\section{Curve representations}

Every curve object contains:
\begin{itemize}[nosep]
  \item a string discriminator \,\field{type};
  \item Cartesian \,\field{start}\, and \,\field{end}\, points; and
  \item the parameters required to reconstruct the analytic path.
\end{itemize}
The endpoints support topology and validation.  The type-specific parameters
define the path between them.  A reader should verify that evaluating those
parameters at $t=0$ and $t=1$ reproduces the supplied endpoints.

\subsection{THETA: constant-colatitude arcs}

Fields are \,\field{theta}, \,\field{phi0}, and signed \,\field{deltaPhi}.
Define
\begin{align}
  \theta(t)&=\theta_*,\\
  \phi(t)&=\phi_0+t\,\Delta\phi.
\end{align}
The Cartesian curve follows from the spherical-to-Cartesian equations in
Section~2.  Its exact length is
\[
  L=R\sin\theta_*\,|\Delta\phi|.
\]
A value $|\Delta\phi|>\pi$ is a legitimate complementary long arc, especially
in polar topology.  Do not normalize it to $(-\pi,\pi]$.

\subsection{PHI: constant-azimuth arcs}

Fields are \,\field{phi}, \,\field{theta0}, and signed
\,\field{deltaTheta}.  Define
\begin{align}
  \theta(t)&=\theta_0+t\,\Delta\theta,\\
  \phi(t)&=\phi_*.
\end{align}
Its exact length is
\[
  L=R|\Delta\theta|.
\]
At a pole, Cartesian coordinates do not determine azimuth.  The exported
\field{phi} value is therefore essential: it specifies the limiting meridian.

\subsection{GENERAL: Cartesian-face/sphere intersection arcs}

A GENERAL curve is an arc of the small circle obtained by intersecting the
sphere with a Cartesian cell-face plane.  Its fields are
\field{center}, \,\field{radius}, \,\field{basisU}, \,\field{basisV},
\field{alpha0}, and signed \,\field{deltaAlpha}.  To avoid confusion with the
sphere radius $R$, denote the JSON curve radius by $\rho$.

Let
\[
  \vect{c}=\field{center},\quad
  \vect{u}=\field{basisU},\quad
  \vect{v}=\field{basisV},\quad
  \alpha(t)=\alpha_0+t\,\Delta\alpha.
\]
Then
\begin{equation}
  \boxed{
  \vect{p}(t)=\vect{c}+\rho\left[
      \vect{u}\cos\alpha(t)+\vect{v}\sin\alpha(t)
  \right].}
  \label{eq:general-curve}
\end{equation}
The vectors $\vect{u}$ and $\vect{v}$ are orthonormal.  Their cross product is
the oriented normal to the circle plane.  The exact derivative is
\begin{equation}
  \vect{p}'(t)=\rho\,\Delta\alpha\left[
      -\vect{u}\sin\alpha(t)+\vect{v}\cos\alpha(t)
  \right],
  \label{eq:general-derivative}
\end{equation}
and the exact arc length is
\[
  L=\rho|\Delta\alpha|.
\]

Useful consistency checks are
\begin{align}
  \|\vect{u}\|&=\|\vect{v}\|=1, &
  \vect{u}\cdot\vect{v}&=0,\\
  \|\vect{p}(t)\|&=R, &
  \vect{p}(0)&=\field{start},\quad
  \vect{p}(1)=\field{end}.
\end{align}

\section{Diagnostics}

Each patch has a \,\field{diagnostics}\, object.

\begin{longtable}{>{\ttfamily}p{0.29\textwidth}p{0.63\textwidth}}
\toprule
\normalfont Field & Meaning \\
\midrule
\endhead
normalizedAreaFraction & Chimera's numerical estimate of $A/(4\pi R^2)$. \\
physicalAreaEstimate & The same estimate multiplied by $4\pi R^2$. \\
perimeter & Sum of exact analytic curve lengths in the declared length unit. \\
enclosesNorthPole & Whether Chimera classifies the patch as enclosing or touching
the north pole. \\
enclosesSouthPole & Analogous south-pole flag. \\
maximumClosureError & Maximum Cartesian distance between the end of one curve
and the start of its successor. \\
\bottomrule
\end{longtable}

The area values are reference diagnostics, not exact answers.  Chimera
currently samples every boundary curve and estimates spherical polygon area by
triangulation.  They are useful for regression comparisons, but an independent
area calculation should be compared against them rather than defined by them.

The root \,\field{summary}\, contains corresponding totals.  In particular,
\field{normalizedAreaFraction} is the sum over final patches and should be near
$1$ for complete coverage.  The field \,\field{intersectingCellCount}\, counts
ordinary non-KISS intersecting cells; \,\field{kissCellCount}\, is separate.
KISS cells are a known special case and do not currently generate ordinary
patches.

\section{Recommended area calculation}

\subsection{Archimedes/Lambert cylindrical equal-area coordinates}

The most useful coordinates for this project are
\begin{equation}
  X=\widetilde{\phi},\qquad Y=\cos\theta=\frac{z}{R},
  \label{eq:equal-area-map}
\end{equation}
where $\widetilde{\phi}$ is a continuously unwrapped azimuth along the directed
boundary.  Since
\[
  dY=-\sin\theta\,d\theta,
\]
the spherical area element satisfies
\[
  dA=R^2\sin\theta\,d\theta\,d\phi
     =-R^2\,dX\,dY.
\]
Therefore the magnitude of spherical patch area equals $R^2$ times the planar
area enclosed in the $(X,Y)$ equal-area plane.

Green's theorem gives either of the equivalent formulas
\begin{align}
  A &= R^2\left|\oint_{\partial P} X\,dY\right|,
  \label{eq:xdy-area}\\
  A &= \frac{R^2}{2}\left|\oint_{\partial P}
          \left(X\,dY-Y\,dX\right)\right|.
  \label{eq:symmetric-area}
\end{align}
Equation~\eqref{eq:xdy-area} is especially attractive here: THETA curves make
zero contribution because $Y$ is constant, PHI curves have a closed-form
contribution, and GENERAL curves require only a smooth one-dimensional
integral.

Do not assume that every exported loop has the same clockwise/counterclockwise
sign in the projected plane.  Compute a signed integral for diagnostics, but
use its absolute value for an individual patch area unless orientation has been
independently established.

\subsection{Curve contributions to $\oint X\,dY$}

For a THETA curve, $Y=\cos\theta_*$ is constant, so
\[
  I_\theta=\int_C X\,dY=0.
\]

For a PHI curve with a continuously selected meridian value
$X=\widetilde{\phi}_*$,
\begin{equation}
  I_\phi=\widetilde{\phi}_*
  \left[\cos(\theta_0+\Delta\theta)-\cos\theta_0\right].
  \label{eq:phi-contribution}
\end{equation}

For a GENERAL curve, let $p_x(t),p_y(t),p_z(t)$ be obtained from
Equation~\eqref{eq:general-curve}.  Then
\begin{align}
  X(t)&=\operatorname{unwrap}\!\left(\operatorname{atan2}
           (p_y(t),p_x(t))\right),\\
  Y(t)&=\frac{p_z(t)}{R},\\
  \frac{dY}{dt}&=\frac{p'_z(t)}{R},
\end{align}
with $p'_z$ available analytically from
Equation~\eqref{eq:general-derivative}.  Its contribution is
\begin{equation}
  I_g=\int_0^1 X(t)\frac{p'_z(t)}{R}\,dt.
  \label{eq:general-contribution}
\end{equation}

Even if a compact symbolic antiderivative for Equation~\eqref{eq:general-contribution}
is inconvenient, evaluating this single analytic boundary integral with
adaptive high-order quadrature is fundamentally different from polygonizing
the entire patch.  It preserves exact curve geometry, provides controlled
error, and is an excellent reference against which a later symbolic result can
be tested.

\subsection{Unwrapping azimuth correctly}

Azimuth unwrapping is the most likely source of a large area error.

\begin{enumerate}
  \item Process curves in their exported boundary order.
  \item Preserve the explicit signed increments on THETA curves.
  \item For sampled or evaluated GENERAL-curve azimuths, add or subtract
        $2\pi$ whenever necessary to maintain continuity with the preceding
        boundary value.
  \item Use the patch's $n_\phi$ interval as an additional consistency check.
        A final patch should occupy one phi band, modulo the identified seam.
  \item At a pole, do not obtain azimuth from \field{atan2(0,0)}.  Use the
        limiting azimuth supplied by the incident PHI or THETA curve.
\end{enumerate}

It is safest to perform the calculation in an \emph{unwrapped} Archimedes
plane.  Do not compute area from a display polygon that has been wrapped,
clipped, or split at the map seam.

\section{Suggested implementation strategy}

The following progression makes failures easier to localize.

\begin{enumerate}
  \item Parse and validate the root format and version.
  \item Validate the grid arrays, sphere radius, and all index ranges.
  \item Reconstruct every curve and verify its two endpoints.
  \item Verify ordered loop closure before attempting any area calculation.
  \item Implement exact THETA and PHI contributions.
  \item Implement Equation~\eqref{eq:general-contribution} using adaptive
        quadrature and explicit azimuth unwrapping.
  \item Compare each result with \,\field{physicalAreaEstimate}\, and record
        absolute and relative differences.  Small patches require an absolute
        tolerance as well as a relative one.
  \item Sum all computed areas and compare with $4\pi R^2$.  Report KISS-cell
        counts alongside this comparison.
  \item Only after the numerical boundary-integral implementation is stable,
        pursue closed-form simplification of the GENERAL contribution if that
        remains a project goal.
\end{enumerate}

Recommended test cases include ordinary nonpolar patches, north- and
south-polar patches, patches adjacent to $\phi=\pm\pi$, boundaries containing a
long THETA arc, very small patches, and reversed synthetic curves for checking
sign handling.

\section{Minimal Java reader and reconstruction example}

Java~17 does not include a JSON parser in the standard library.  The following
example uses Jackson, the same library used by Chimera.  Add this dependency to
the student's Maven project (or use a compatible newer Jackson 2.x release).

\begin{lstlisting}[language=XML,caption={Maven dependency for the example}]
<dependency>
  <groupId>com.fasterxml.jackson.core</groupId>
  <artifactId>jackson-databind</artifactId>
  <version>2.17.2</version>
</dependency>
\end{lstlisting}

Listing~\ref{lst:java-reader} reads the document as a Jackson tree.  This keeps
the example compact while retaining explicit validation and type dispatch.  A
larger application may replace the tree with records or classes after the file
format is understood.  Production code should additionally validate against
the supplied JSON Schema.

\begin{lstlisting}[language=Java,caption={Java 17 parser, curve evaluator, and closure check},label={lst:java-reader}]
import java.io.IOException;
import java.nio.file.Path;

import com.fasterxml.jackson.databind.JsonNode;
import com.fasterxml.jackson.databind.ObjectMapper;

public final class ChimeraPatchReader {

    private static final double TOLERANCE = 1.0e-7;
    private static final ObjectMapper JSON = new ObjectMapper();

    private ChimeraPatchReader() { }

    public static void main(String[] args) throws IOException {
        Path input = Path.of(args.length == 0
                ? "Paper-Test-Grid-patches.json" : args[0]);
        JsonNode document = JSON.readTree(input.toFile());

        require("mdi-chimera-patches".equals(
                document.path("format").asText()),
                "Not an MDI Chimera patch file");
        require(document.path("version").asInt(-1) == 1,
                "Unsupported Chimera patch version");

        double sphereRadius = document.path("grid").path("sphere")
                .path("radius").asDouble(Double.NaN);
        require(Double.isFinite(sphereRadius) && sphereRadius > 0.0,
                "Invalid sphere radius");

        JsonNode patches = document.path("patches");
        require(patches.isArray(), "patches must be an array");

        for (JsonNode patch : patches) {
            validatePatch(patch, sphereRadius);
        }

        System.out.printf("Read and validated %,d final patches.%n",
                patches.size());
    }

    private static void validatePatch(JsonNode patch, double sphereRadius) {
        int patchId = patch.path("id").asInt(-1);
        JsonNode boundary = patch.path("boundary");
        require(boundary.isArray() && boundary.size() >= 2,
                "Patch " + patchId + " has an incomplete boundary");

        for (int i = 0; i < boundary.size(); i++) {
            JsonNode curve = boundary.get(i);
            JsonNode next = boundary.get((i + 1) % boundary.size());

            double[] exportedStart = vector(curve, "start");
            double[] exportedEnd = vector(curve, "end");
            double[] nextStart = vector(next, "start");
            double[] reconstructedStart = evaluate(curve, 0.0,
                    sphereRadius);
            double[] reconstructedEnd = evaluate(curve, 1.0,
                    sphereRadius);

            require(distance(exportedStart, reconstructedStart) <= TOLERANCE,
                    "Curve parameters do not reproduce the start point");
            require(distance(exportedEnd, reconstructedEnd) <= TOLERANCE,
                    "Curve parameters do not reproduce the end point");
            require(distance(exportedEnd, nextStart) <= TOLERANCE,
                    "Patch " + patchId + " is open at curve " + i);
            require(Math.abs(norm(exportedStart) - sphereRadius)
                    <= TOLERANCE, "Curve start is not on the sphere");
        }
    }

    public static double[] evaluate(JsonNode curve, double t,
            double sphereRadius) {
        require(t >= 0.0 && t <= 1.0, "t is outside [0,1]");
        String type = curve.path("type").asText();

        return switch (type) {
        case "THETA" -> {
            double theta = number(curve, "theta");
            double phi = number(curve, "phi0")
                    + t * number(curve, "deltaPhi");
            yield sphericalPoint(sphereRadius, theta, phi);
        }
        case "PHI" -> {
            double theta = number(curve, "theta0")
                    + t * number(curve, "deltaTheta");
            yield sphericalPoint(sphereRadius, theta,
                    number(curve, "phi"));
        }
        case "GENERAL" -> {
            double[] center = vector(curve, "center");
            double[] u = vector(curve, "basisU");
            double[] v = vector(curve, "basisV");
            double rho = number(curve, "radius");
            double alpha = number(curve, "alpha0")
                    + t * number(curve, "deltaAlpha");
            double cosine = Math.cos(alpha);
            double sine = Math.sin(alpha);
            yield new double[] {
                center[0] + rho * (u[0] * cosine + v[0] * sine),
                center[1] + rho * (u[1] * cosine + v[1] * sine),
                center[2] + rho * (u[2] * cosine + v[2] * sine)
            };
        }
        default -> throw new IllegalArgumentException(
                "Unknown curve type: " + type);
        };
    }

    private static double[] sphericalPoint(double radius, double theta,
            double phi) {
        double sinTheta = Math.sin(theta);
        return new double[] {
            radius * sinTheta * Math.cos(phi),
            radius * sinTheta * Math.sin(phi),
            radius * Math.cos(theta)
        };
    }

    private static double number(JsonNode object, String field) {
        double value = object.path(field).asDouble(Double.NaN);
        require(Double.isFinite(value), field + " is missing or non-finite");
        return value;
    }

    private static double[] vector(JsonNode object, String field) {
        JsonNode value = object.path(field);
        require(value.isArray() && value.size() == 3,
                field + " must contain three numbers");
        return new double[] {
            finite(value.get(0).asDouble(Double.NaN), field),
            finite(value.get(1).asDouble(Double.NaN), field),
            finite(value.get(2).asDouble(Double.NaN), field)
        };
    }

    private static double finite(double value, String label) {
        require(Double.isFinite(value), label + " contains a non-finite value");
        return value;
    }

    private static double norm(double[] value) {
        return Math.sqrt(value[0] * value[0]
                + value[1] * value[1] + value[2] * value[2]);
    }

    private static double distance(double[] a, double[] b) {
        double dx = a[0] - b[0];
        double dy = a[1] - b[1];
        double dz = a[2] - b[2];
        return Math.sqrt(dx * dx + dy * dy + dz * dz);
    }

    private static void require(boolean condition, String message) {
        if (!condition) {
            throw new IllegalArgumentException(message);
        }
    }
}
\end{lstlisting}

The example deliberately preserves every signed angular increment.  The next
step for the student is to retain the parsed boundary order and apply
Equations~\eqref{eq:phi-contribution} and
\eqref{eq:general-contribution} to those same curve objects.  In a production
reader, also validate five-tuple ranges and uniqueness, GENERAL-curve basis
orthonormality, and the complete grid arrays as described below.

\section{Validation checklist}

Before trusting computed areas, verify all of the following.

\begin{itemize}[label=$\square$]
  \item Root format is \,\field{mdi-chimera-patches}\, and version is $1$.
  \item All required numerical values are finite.
  \item Grid arrays are strictly increasing and span their declared domains.
  \item The sphere is centered at the origin and has positive radius.
  \item Every five-tuple is in range and unique.
  \item Every boundary contains at least two curves.
  \item Every analytic curve reproduces its exported start and end.
  \item GENERAL bases are unit length and mutually orthogonal.
  \item Every endpoint lies on the sphere within tolerance.
  \item Every curve end matches the next curve start, including loop closure.
  \item Signed angular sweeps are preserved without unwanted normalization.
  \item Azimuth is continuous after unwrapping, including seam and pole cases.
  \item Per-patch areas are nonnegative and finite.
  \item The sum of areas is compared with both Chimera's exported total and
        $4\pi R^2$.
\end{itemize}

\section{Interpretation of discrepancies}

A discrepancy localized to seam-adjacent patches usually indicates incorrect
azimuth unwrapping.  A discrepancy localized to polar patches usually indicates
that azimuth was inferred from a pole's Cartesian coordinates rather than from
the incident meridian.  Errors only on long THETA arcs usually indicate that a
signed sweep was normalized to its short complement.  Widespread small errors
that shrink with tighter quadrature tolerance suggest integration error;
widespread errors that do not shrink suggest a convention or scale error.

For each comparison, retain the patch ID, five-tuple, signed integral, absolute
area, Chimera estimate, absolute difference, relative difference, pole flags,
and seam-adjacency status.  That diagnostic table will make the eventual report
far more useful than a single global error number.

\section{Final advice}

Treat topology and integration as separate problems.  First prove that the
ordered directed loop has been reconstructed correctly; only then integrate it.
Start with a controlled numerical evaluation of the exact boundary formulas,
because it provides a high-quality benchmark for any later symbolic derivation.
Use Chimera's diagnostics as independent comparisons, not as values to copy.
Finally, preserve failing patch IDs and five-tuples: the difficult seam, pole,
and near-degenerate cases are the most scientifically valuable tests of the
method.

\end{document}
